% compendium/execution-trading-costs/execution-benchmarks.tex
% Entries migrated on 2026-10-15 from the legacy execution-benchmarks file (slice 11), owned by the
% section overview execution-trading-costs; dispositions and source checks in notes/migration.md,
% the blocks and the original table verbatim in notes/artefacts/market-making-lob/derivations.tex.

\section{Execution benchmarks and implementation shortfall}

\begin{qfdefinition}[Execution benchmark: a reference price for judging an execution, or the trading target of an algorithm]\label{def:execution-benchmark}
\qffield{Statement} The word has two uses.
\begin{itemize}
\item A \emph{benchmark price} $P_B$ against which realised execution is measured. With average execution price $\bar P$ and $\mathrm{Side}=+1$ for a buy and $-1$ for a sell, the cost in basis points is $\mathrm{Side}\,(\bar P-P_B)/P_B\times10^4$, positive when the execution was worse than the benchmark; the profit-and-loss form multiplies it by $-1$ \citep[chap.~3, eqs.~(3.11)--(3.14), pp.~105--106]{s7-kissellscience}. Common benchmark prices are the quotation midpoint when the manager decided to trade (Perold's implementation shortfall, \hyperref[def:implementation-shortfall]{below}), the price when the order arrived in the market, the volume-weighted average price (VWAP) of the day or of the order's own trading interval, a participation-weighted price, and the open, the close, the next day's open or close, or the average of the day's open, high, low and close \citep[chap.~``Liquidity and Transaction Cost Measurement'', pp.~424--426]{alpha-harris}; \citep[chap.~3, pp.~106--108]{s7-kissellscience}.
\item The \emph{target} of an algorithm class. A VWAP algorithm trades in proportion to the intraday volume curve; a TWAP algorithm trades equal amounts in equal time intervals, about 1/390 of a full-day US order per minute, often skipping the opening and closing auctions; a percentage-of-volume (POV) algorithm trades a fixed fraction of market volume until the order is done; arrival-price and implementation-shortfall algorithms balance expected cost against risk at a chosen risk aversion, the latter adding real-time adaptation, and their meaning differs across brokers \citep[chap.~1, pp.~16--18]{s7-kissellscience}. A liquidity seeker has no reference price: it is relative-price oriented, captures liquidity across pools and minimises information leakage \citep[sec.~3.4.1, Table~3.2, p.~230]{alpha-laruelle}.
\end{itemize}

\qffield{Notation} What each benchmark price fixes, from the passages checked (the order's own fills are $s_j$ shares at $p_j$; the unfilled remainder is treated in \hyperref[def:implementation-shortfall]{implementation shortfall}):
\begin{itemize}
\item \emph{Decision price}: the quotation midpoint when the manager decided to trade \citep[chap.~``Liquidity and Transaction Cost Measurement'', p.~426]{alpha-harris}; in practice most managers use the previous day's close \citep[same chap., box ``The Plexus Iceberg of Transaction Costs'', p.~437]{alpha-harris}. Known before the order is worked.
\item \emph{Arrival price} $P_0$: the price when the order is entered in the market, the bid--ask midpoint in market hours, the last trade or official close otherwise \citep[chap.~3, p.~102]{s7-kissellscience}. Known when trading starts.
\item \emph{VWAP}: total traded value over total volume, each trade weighted by its size \citep[chap.~``Liquidity and Transaction Cost Measurement'', p.~425]{alpha-harris}. Full-day VWAP has no official value; vendors choose the eligible trades and may drop delayed or negotiated ones. Interval VWAP covers the order's own trading horizon, computed from tick data; VWAP to the end of the day judges orders finished early \citep[chap.~3, pp.~107--108]{s7-kissellscience}. Known only afterwards, and it includes the order's own trades.
\item \emph{Participation-weighted price} at rate $\rho$: the VWAP of the first (shares traded)$/\rho$ shares of market volume from the arrival time \citep[chap.~3, eqs.~(3.15)--(3.16), p.~108]{s7-kissellscience}. Known only afterwards.
\item \emph{Open, close, later closes}: single daily prices; next-day and later prices are used to separate temporary from permanent impact \citep[chap.~3, p.~106]{s7-kissellscience}.
\end{itemize}
Kissell's benchmark formulas use the average execution price without fees, so they are gross of explicit costs; his implementation shortfall adds fees \citep[chap.~3, eqs.~(3.4)--(3.16), pp.~98--108]{s7-kissellscience}. Signs are not standard across the industry: some report a cost, positive when worse, others a gain, positive when better \citep[chap.~3, pp.~97 and~108]{s7-kissellscience}.

\qffield{Conditions} Harris compares the benchmark prices on four properties \citep[chap.~``Liquidity and Transaction Cost Measurement'', sec.~21.4, pp.~427--432, Table~21-1, p.~433]{alpha-harris}:
\begin{itemize}
\item \emph{noise}, which grows with the time between trade and benchmark: the contemporaneous midpoint is least noisy, open and close most;
\item \emph{split orders}: a benchmark moved by the order's own impact understates its cost. A trader who is the only buyer of the day matches VWAP whatever the impact (pp.~428--429);
\item \emph{momentum and contrarian trading}, which bias open and, less, VWAP benchmarks; \emph{informed trading}, which makes the close and, less, VWAP understate costs;
\item \emph{gaming}: a broker with discretion over timing can game any benchmark that depends on when it trades, for example by spreading fills through the day to match VWAP or trading at the close. A benchmark fixed before the broker receives the order cannot be gamed, so implementation shortfall cannot (pp.~430--431).
\end{itemize}

\qffield{Practice} When data cost is no object, implementation shortfall is the best estimator; retail traders use the effective spread, and sponsors who receive only daily average prices use VWAP \citep[same chap., pp.~426 and~431--432]{alpha-harris}. VWAP results are not comparable across stocks of different volatility, and converge to zero as the order becomes a large share of volume \citep[chap.~3, p.~107]{s7-kissellscience}. Single-trade estimates are noisy; averaging reduces noise but not bias or gaming, and comparisons across brokers are meaningful only for similar trading problems \citep[chap.~``Liquidity and Transaction Cost Measurement'', p.~432]{alpha-harris}.

\qffield{Pitfalls} Interpretation: a benchmark price is a reference, not an executable price. A midpoint benchmark charges an aggressive order at least the half-spread, as the \hyperref[def:spread-measures]{effective spread} does. Decision and arrival prices are known when the order is planned; VWAP, participation-weighted and closing prices are known only after it, so an algorithm can target them only through forecasts, such as a forecast volume curve. Performance against the benchmark an algorithm targets is not its cost: a VWAP-tracking execution can match VWAP and still have paid spread and impact against arrival. Kissell calls TWAP ``a constant participation rate'' while describing equal shares per minute, which is a constant rate in time, not in volume \citep[chap.~1, p.~17]{s7-kissellscience}. His participation-weighted formula (3.16) has no side factor and is correct as written only for buys. Lehalle and Laruelle's ``benchmarks'' in Table~3.1 are algorithm classes; PoV and liquidity seekers are not benchmark prices (see \hyperref[meth:algorithm-benchmark-choice]{benchmark choice}). No passage checked defines a TWAP benchmark price.

\qffield{See also} \hyperref[def:implementation-shortfall]{implementation shortfall} (the decision-price benchmark with its unfilled remainder); \hyperref[meth:algorithm-benchmark-choice]{choosing an algorithm's benchmark}; \hyperref[def:spread-measures]{spread measures} (contemporaneous and post-trade midpoint benchmarks); \hyperref[ev:intraday-impact-profile]{intraday impact profile} (impact measured from the arrival price); \hyperref[mod:almgren-chriss-execution]{mean--variance execution} (the schedule behind arrival-price algorithms).

\qffield{Further reading} The legacy table it replaces in part, verbatim: \nolinkurl{notes/artefacts/market-making-lob/derivations.tex}.

\qffield{Index} execution benchmark; benchmark price; decision price; arrival price; VWAP; interval VWAP; full-day VWAP; VWAP to end of day; TWAP; participation-weighted price; PWP; percentage of volume; POV; PoV; opening price; closing price; transaction cost analysis; TCA; split orders; gaming; benchmark noise; liquidity seeker

\qffield{Status} Verification: partly checked 2026-10-15 (against the text layers of local copies of \texttt{alpha-harris}, 2002 Oxford printing, pp.~424--437, Table~21-1 on p.~433 by page image; \texttt{s7-kissellscience}, Academic Press printing with copyright 2014, ISBN 978-0-12-401689-7, chap.~1, pp.~16--18, and chap.~3, pp.~97--108, pp.~100, 103 and~105 by page image; and \texttt{alpha-laruelle}, 2018 World Scientific printing, p.~230). Stability: draft. Currency: definitions timeless; vendor and consultant practice as of 2002 and 2014.
\end{qfdefinition}

\begin{qfdefinition}[Implementation shortfall: paper portfolio minus real portfolio, with execution, delay and opportunity components]\label{def:implementation-shortfall}
\qffield{Statement} Implementation shortfall is the paper return of a portfolio assumed to trade the whole order at the manager's decision price, minus the actual return, using actual fills and fees \citep[chap.~3, eqs.~(3.1)--(3.4), pp.~97--99]{s7-kissellscience}. For an intended order of $S$ shares, fills of $s_j$ shares at $p_j$, decision price $P_d$, price $P_n$ at the end of the measurement period and fees,
\[ \mathrm{IS}=\underbrace{\textstyle\sum_js_jp_j-\sum_js_jP_d}_{\text{execution}}+\underbrace{\bigl(S-\textstyle\sum_js_j\bigr)(P_n-P_d)}_{\text{opportunity}}+\text{fees}, \]
which Kissell calls Perold's formulation \citep[same chap., pp.~100--101, eq.~(3.6), p.~104]{s7-kissellscience}. With the arrival price $P_0$, Wagner's expanded form splits it into delay, trading and opportunity parts:
\[ \mathrm{IS}=S(P_0-P_d)+\textstyle\sum_js_j\,(\bar P-P_0)+\bigl(S-\sum_js_j\bigr)(P_n-P_0)+\text{fees},\qquad\bar P=\sum_js_jp_j\big/\sum_js_j \]
\citep[same chap., pp.~101--104, eq.~(3.7)]{s7-kissellscience}. The trading-related part is the arrival cost \citep[eq.~(3.9), p.~105]{s7-kissellscience}. Harris states the same estimator with the decision-time midpoint: the filled size times the signed gap between average fill and benchmark, plus the unfilled size times the gap between a later price and the benchmark \citep[chap.~``Liquidity and Transaction Cost Measurement'', p.~426]{alpha-harris}.

\qffield{Notation} Quantities are signed, $S,s_j>0$ for buys and $<0$ for sells \citep[chap.~3, eq.~(3.9), p.~105]{s7-kissellscience}; IS is in currency and positive when the real portfolio did worse. Fees are explicit costs: commissions, taxes, clearing and settlement charges, ticket charges and rebates (p.~98). The spread and impact are inside the fill prices.

\qffield{Conditions} When the order fills completely, IS reduces to $S(\bar P-P_d)+\text{fees}$ and does not depend on $P_n$ (pp.~99--100). An unfilled or cancelled remainder is charged at $P_n$, so its cost depends on when it is measured: at the end of trading \citep[chap.~3, p.~102]{s7-kissellscience}, or a day, a week or a month later; the Plexus Group uses 30 trading days after the decision \citep[chap.~``Liquidity and Transaction Cost Measurement'', pp.~436--437]{alpha-harris}. The remainder counts only for an order that could realistically have been filled: Harris's 4-million-share sell order, 50 days' volume, would show a missed-trade cost of \$714{,}600 that the trader never had the opportunity to avoid (p.~436).

\qffield{Practice} The split shows who is responsible: delay belongs to the manager or desk, trading-related cost to the broker or algorithm, and the arrival cost is the main measure of broker and algorithm performance \citep[chap.~3, pp.~101--104]{s7-kissellscience}. When the decision price is unknown, only the market-activity part from $P_0$ is measured (p.~104, eq.~(3.8)). Plexus's components, manager timing, trader timing, market impact, commission and missed trades, sum to the total; Plexus estimated its clients' costs at 12 basis points of commission, 20 of impact, 53 of timing and 16 of missed trades \citep[chap.~``Liquidity and Transaction Cost Measurement'', box, p.~437]{alpha-harris}. Interpretation, checked with exact arithmetic: a buy of 10{,}000 shares decided at a midpoint of 50.00, arriving at 50.10, filling 3{,}000 at 50.20 and 3{,}000 at 50.30, cancelling the rest, with $P_n=50.60$ and fees of 60, has paper gain 6{,}000 and real gain 2{,}040, so $\mathrm{IS}=3{,}960$: Perold $1{,}500+2{,}400+60$; Wagner delay $1{,}000$, trading $900$, opportunity $2{,}000$, fees $60$. The mirror-image sell gives the same figures with signed quantities, and Kissell's worked examples (\$2{,}600, and \$3{,}080 split as 1{,}250, 1{,}000, 750 and 80) reconcile.

\qffield{Pitfalls} Interpretation: Kissell's intermediate actual return for a partial fill adds the fees instead of subtracting them (p.~100); his final formula is right. His basis-point arrival cost (3.10) divides by the signed $(\sum s_j)P_0$, which reverses the sign for sells; (3.11), with the side factor, does not (in the sell above, $-30.06$ against $+30.06$ basis points). Count each cost once: with a midpoint decision or arrival price, the half-spread is already inside the execution component, so adding an effective-spread term double-counts it, and fees enter only as the explicit term. Splitting the trading-related part into spread and impact, as the overview's checkpoint asks, needs a further benchmark such as the effective spread at each fill; none of the passages checked does it. The \hyperref[mod:almgren-chriss-execution]{Almgren--Chriss} shortfall $XS_0-\sum n_k\tilde S_k$ is the trading-related part for a complete sale measured from the pre-trade price, with half-spread and fees inside the temporary impact; its mean and variance are model forecasts, not a realised shortfall, and it has no opportunity term because the model forces completion. Perold's paper itself (\texttt{wide-perold-shortfall}) was not available: the attribution of the formula to Perold, and of the expanded form to Wagner (1975, 1991), is Kissell's; Harris credits Perold with popularising the decision-midpoint estimator and Treynor, writing seven years earlier, with the paper-portfolio method (p.~424).

\qffield{See also} \hyperref[def:execution-benchmark]{execution benchmarks} (the other benchmark prices); \hyperref[mod:almgren-chriss-execution]{mean--variance execution} (expected shortfall and its variance); \hyperref[def:spread-measures]{spread measures}; \hyperref[def:price-impact]{price impact}.

\qffield{Further reading} \texttt{wide-perold-shortfall}, the overview's first reading, not yet read. The derivations artefact: \nolinkurl{notes/artefacts/market-making-lob/derivations.tex}.

\qffield{Index} implementation shortfall; IS; paper portfolio; paper return; decision price; arrival price; arrival cost; execution cost; delay cost; trading-related cost; opportunity cost; missed trade opportunity cost; expanded implementation shortfall; Wagner; Perold; Treynor; Plexus; fees; commissions; partial fill; unfilled order; cancellation

\qffield{Status} Verification: partly checked 2026-10-15 (against the text layers of local copies of \texttt{alpha-harris}, 2002 Oxford printing, pp.~424, 426 and~434--437, and \texttt{s7-kissellscience}, Academic Press printing with copyright 2014, chap.~3, pp.~97--105, pp.~100, 103 and~105 by page image; formulas and examples checked with exact arithmetic; \texttt{wide-perold-shortfall} not consulted). Stability: draft. Currency: timeless (accounting); the Plexus figures as of 2002.
\end{qfdefinition}

\begin{qfmethod}[Choosing an algorithm's benchmark from order size, market context and the risk to hedge (Lehalle and Laruelle)]\label{meth:algorithm-benchmark-choice}
\qffield{Statement} Lehalle and Laruelle tie each benchmark of a trading algorithm to the investor needs in a region: in the US, where high-frequency traders are everywhere, stealth; in Asia, comparing trading efficiency across local markets, given their differences; in Europe, both avoiding detection by high-frequency traders and learning how efficiently a set of order books is reached, to adjust routing quickly. ``On paper'' European investors should therefore target implementation shortfall, Asian ones VWAP or PoV, US ones liquidity seekers. Size and market conditions adjust these ``drastic rules'': liquidity seekers suit orders without minimum participation constraints, IS small orders, VWAP larger and PoV very large ones \citep[sec.~3.4.1, pp.~227--228]{alpha-laruelle}. Their Table~3.1 summarises the remarks (p.~229):

\smallskip
\noindent{\small\begin{tabularx}{\linewidth}{@{}>{\RaggedRight}p{0.2\linewidth}>{\RaggedRight}p{0.13\linewidth}>{\RaggedRight}p{0.17\linewidth}>{\RaggedRight}p{0.17\linewidth}>{\RaggedRight\arraybackslash}X@{}}
\toprule
Benchmark & Region & Order size & Market context & Risk hedged (do not miss)\\
\midrule
PoV (\hyperref[def:execution-benchmark]{target}) & Asia & more than 10\% of ADV & possible negative news & the rapid spread of unexpected news, especially if the investor has the information\\
VWAP/TWAP (\hyperref[def:execution-benchmark]{target}) & Asia, Europe & 5--15\% of ADV & unusual volume negligible & the slow spread of information\\
Implementation shortfall (\hyperref[def:implementation-shortfall]{measure}) & Europe, US & 0--6\% of ADV & possible price opportunities & an unexpected price move\\
Liquidity seeker & US & any & price oscillates around fair value & a liquidity burst or relative price move\\
\bottomrule
\end{tabularx}}
\smallskip

\noindent ADV is average daily consolidated volume. Sources: rows from \citep[sec.~3.4.1, Table~3.1, p.~229]{alpha-laruelle}, checked by page image; links to this compendium's definitions. The table is a summary of the entries, and the risk column is condensed from the source's wording.

\qffield{Conditions} Lehalle and Laruelle's guidance is practitioner judgement: the passage gives no data for it, and the table ``should not be read as `all conditions in the columns have to be fulfilled''', only as possible reasons for a benchmark (p.~228). The size bands overlap as printed (0--6\%, 5--15\%, above 10\%).

\qffield{Practice} The hedged-risk column matters because an algorithm performs a risk--reward arbitrage: once its risk is under control it is free to adjust the prices and quantities of its orders, so optimal execution combines risk control, by following an ``optimal'' trading rate, with capturing price and liquidity opportunities. The choice also depends on the stock's liquidity and on the investment style behind the trade: hedging trades are more exposed to execution-specific impact, long-term position building to ``systemic'' impact that algorithms cannot do much about \citep[sec.~3.4.1, p.~228]{alpha-laruelle}. Table~3.2 adds the type of stock, trade and main feature of each algorithm (p.~230).

\qffield{Pitfalls} The legacy notes said VWAP and PoV help Asia compare ``fragmented market conventions''; the source speaks of differences between local markets, not fragmentation. They gave Europe only routing discipline, omitting avoidance of detection by high-frequency traders. Their risk column dropped ``especially if I have the information'' for PoV. Interpretation: the legacy caveat that a benchmark ``justifies the comparison, not the whole execution policy'' is consistent with the risk--reward reading on p.~228 but is not in the source. The algorithm a trade is given does not settle how its cost is measured; an order worked by a VWAP algorithm can still be judged against arrival price, and matching VWAP does not show low cost (\hyperref[def:execution-benchmark]{execution benchmarks}). The regional pattern is dated by the 2018 printing.

\qffield{See also} \hyperref[def:execution-benchmark]{execution benchmarks}; \hyperref[def:implementation-shortfall]{implementation shortfall}; \hyperref[mod:almgren-chriss-execution]{mean--variance execution} (the risk-control layer and its trading curve); \hyperref[def:block-trade]{block trades}. Fragmentation, dark trading and venue rules are the subject of the specialist pathway \texttt{market-structure-execution-quality}, which has no entries yet.

\qffield{Further reading} The legacy blocks and table, verbatim: \nolinkurl{notes/artefacts/market-making-lob/derivations.tex}.

\qffield{Index} benchmark choice; algorithm selection; trading algorithm benchmark; PoV; percentage of volume; VWAP; TWAP; implementation shortfall; liquidity seeker; regional benchmark preference; ADV; hedged risk; stealth trading; Lehalle--Laruelle Table 3.1

\qffield{Status} Verification: partly checked 2026-10-15 (against the text layer of a local copy of \texttt{alpha-laruelle}, 2018 World Scientific printing, sec.~3.4.1, pp.~227--230, Table~3.1 on p.~229 by page image). Stability: draft. Currency: practitioner guidance as of the 2018 printing.
\end{qfmethod}
